\documentclass[a4paper,oneside,onecolumn,12pt]{LegrandOrangeBook}
\usepackage[utf8]{inputenc}
\usepackage[T1]{fontenc}
\def\magyarOptions{chapterhead=unchanged}
\usepackage[magyar]{babel}
\usepackage{longtable}
\usepackage{setspace}
\usepackage{parskip}
\usepackage{minted}
\usepackage{svg}
\usepackage{ragged2e}
\usepackage{markdown}
\usepackage[SC5b,ISBN=000-00-000-0000-0]{ean13isbn}

\markdownSetup{pipeTables,tableCaptions}

\newcommand{\chpt}[1]{\chapter*{#1}\addcontentsline{toc}{section}{#1}}

% TODO: finish?
\hypersetup{
    pdftitle={A szakdolgozat címe},            % Title field
    pdfauthor={Szerző Neve},                   % Author field
    pdfsubject={Szakdolgozat v. Diplomamunka}, % Subject field
    pdfkeywords={key1, key2, key3, key4},      % Keywords
    pdfcreator={LaTeX},                        % Content creator field
}

% Define the color used for highlighting throughout the book
\definecolor{ocre}{RGB}{243, 102, 25}
% Default whitespace from the top of the page to the chapter title on chapter pages
\chapterspaceabove{6.5cm}
% Default amount of vertical whitespace from the top margin to the start of the text on chapter pages
\chapterspacebelow{6.75cm}

\onehalfspacing
\frenchspacing

\hypersetup{
    % pdftitle={A szakdolgozat címe}
    colorlinks=true,
    linkcolor=blue,
    filecolor=magenta,      
    urlcolor=cyan,
}

\urlstyle{same}

\ExecuteBibliographyOptions{sorting=none}
\addbibresource{manuscript.bib}

\begin{document}

% Title page
\thispagestyle{empty}
% TODO: finish?
\begin{minipage}[c][\textheight][c]{\textwidth}
    {
        \centering
        \includegraphics[keepaspectratio,width=3cm]{SelyeBanner.png} \\
        \vskip0.5cm
        {\LARGE UNIVERZITA J. SELYEHO} \\
        \vskip0.5cm
        {\LARGE SELYE JÁNOS EGYETEM} \\
        \vskip0.5cm
        {\large Fakulta ekonómie a informatiky} \\
        \vskip0.5cm
        {\large Gazdaságtudományi és Informatikai Kar} \\
        \vfill
        {\Huge A szakdolgozat címe} \\
        Szakdolgozat v. Diplomamunka \\
        Szerző Neve \\
        \ISBN \\
        \hfill\the\year{}, Komárom\hfill
    }
\end{minipage}
\cleardoublepage

\begingroup
\makeatletter
\let\ps@plain\ps@empty
% TODO: finish?
\begin{minipage}[c][\textheight][c]{\textwidth}
	{
        \centering
	    {\large UNIVERZITA J. SELYEHO\\SELYE JÁNOS EGYETEM} \\
	    \vskip0.5cm
        {\ NÁZOV FAKULTY\\Fakulta ekonómie a informatiky\\Gazdaságtudományi és Informatikai Kar} \\
        \vfill
        {\Large NÁZOV PRÁCE\\Vytvorenie softvérového interpretera pre CHIP-8 } \\
        \vfill
        \thispagestyle{empty}
        \begin{tabular}{ll}
            Študijný program:    & Aplikovaná informatika \\
            Tanulmányi program:  & Alkalmazott Informatika \\
            Študijný odbor:      & Informatika \\
            Tanulmányi szak:     & Informatika \\
            Školiteľ:            & Tanár Neve \\
            Témavezető:          & Tanár Neve \\
            Konzultant:          & Tanár Neve \\
            Konzulens:           & Tanár Neve \\
            Školiace pracovisko: & Katedra informatiky \\
            Tanszék megnevezése: & Informatikai Tanszék \\
        \end{tabular}
        \vfill
        Označenie typu práce - Szakdolgozat v. Diplomamunka \\
        Szerző Neve \\
        \ISBN \\
        \hfill\the\year{}, Komárno\hfill
	}
\end{minipage}
\endgroup

{
    \hspace*{-2cm}
    % \includegraphics[keepaspectratio, width=17cm]{./zadanie-zp_17422.pdf}
    Ide jön az aláírt témakiírás
}

\tableofcontents
\pagebreak
\listoffigures
\addcontentsline{toc}{section}{Ábrák jegyzéke}
\pagebreak

% The begining of the contents
% \chapterimage{kep/header.png}

\chpt{Feladatkiírás}
A dolgozat célja egy saját programnyelv készítése a hozzá tartozó fordítóprogrammal és mintaprogramokkal. A szerző egy saját programnyelvet tervez (Benex). Ehhez a programnyelvhez készít egy fordító programot, mely tartalmaz egy frontendet (parser) és egy backendet is (compiler). A compiler futtatható fájlokat hoz létre. A szerző demonstrálja az elkészített compilert néhány példaprogram segítségével.

% \chapterimage{kep/header2.png}
\chpt{Opis práce}
Cieľom tejto práce je vytvoriť programovací jazyk s kompilátorom a vzorovými programami. Autor vyvinie vlastný programovací jazyk (Benex). Pre tento programovací jazyk autor vyvíja kompilátor, ktorý obsahuje frontend (parser) a backend (kompilátor). Kompilátor vytvára spustiteľné súbory. Autor demonštruje kompilátor na niekoľkých príkladoch programov.

\pagebreak

% \chapterimage{kep/header3.png}
\chpt{Abstrakt}\label{sec:abstrakt}
Cieľom tejto práce je vytvoriť programovací jazyk s kompilátorom a vzorovými programami. Autor vyvinie vlastný programovací jazyk (Benex). Pre tento programovací jazyk autor vyvíja kompilátor, ktorý obsahuje frontend (parser) a backend (kompilátor). Kompilátor vytvára spustiteľné súbory. Autor demonštruje kompilátor na niekoľkých príkladoch programov.

\textbf{Kľúčové slová: klúč1, klúč2, klúč3}

\pagebreak

% \chapterimage{kep/header4.png}
\chpt{Absztrakt}\label{sec:absztrakt}
A dolgozat célja egy saját programnyelv készítése a hozzá tartozó fordítóprogrammal és mintaprogramokkal. A szerző egy saját programnyelvet tervez (Benex). Ehhez a programnyelvhez készít egy fordító programot, mely tartalmaz egy frontendet (parser) és egy backendet is (compiler). A compiler futtatható fájlokat hoz létre. A szerző demonstrálja az elkészített compilert néhány példaprogram segítségével.

\textbf{Kulcsszavak: kulcs1, kulcs2, kulcs3}

\pagebreak

% \chapterimage{kep/header.png}
\chpt{Abstract}
The aim of this thesis is to create a programming language with compiler and sample programs. The author develops an own programming language (Benex). For this programming language, the author develops a compiler, which includes a frontend (parser) and a backend (compiler). The compiler produces executable files. The author demonstrates the compiler with some example programs.

\textbf{Keywords: key1, key2, key3}

% TODO: intro here
\chapter*{Bevezetés}
\addcontentsline{toc}{chapter}{Bevezetés}
\markboth{}{\sffamily\normalsize{Bevezetés}}
\begin{eBox}
    This is a comment that i might want to add ?? i guess
\end{eBox}
This is the place where i write about the intro...

% Begining of the main content
\chapter{Language Theory}
\markdownInput{notes.md}
\markdownInput{x86.md}

% \chapterimage{kep/header2.png}
\chapter{The Compiler}
\markdownInput{structure.md}
\markdownInput{spec.md}
\markdownInput{functions.md}

% \chapterimage{kep/header.png}
\chapter*{Befejezés}
\addcontentsline{toc}{chapter}{Befejezés}
Some conclusion of the work.

\pagebreak

% \chapterimage{kep/header3.png}
\chapter*{Resumé}
\addcontentsline{toc}{chapter}{Resumé}
\markboth{}{\sffamily\normalsize{Resumé}}
Some final resume of the work.

% TODO: add this once there is something in the bibliography
% \printbibliography[title=Hivatkozások]
% \addcontentsline{toc}{chapter}{Hivatkozások}

\pagebreak
\thispagestyle{empty}

\mbox{}
\vfill

% \begin{Center}
%     \mbox{\vskip1cm}
%     \EANisbn
% \end{Center}

\end{document}
